%% LyX 2.4.4 created this file.  For more info, see https://www.lyx.org/.
%% Do not edit unless you really know what you are doing.
\documentclass[english]{article}
\usepackage[T1]{fontenc}
\usepackage[latin9]{inputenc}
\usepackage{units}
\usepackage{enumitem}
\usepackage{amsmath}
\usepackage{amsthm}
\usepackage{graphicx}
\usepackage{geometry}
\geometry{verbose,tmargin=1in,bmargin=1in,lmargin=1in,rmargin=1in}
\PassOptionsToPackage{normalem}{ulem}
\usepackage{ulem}

\makeatletter
%%%%%%%%%%%%%%%%%%%%%%%%%%%%%% Textclass specific LaTeX commands.
\numberwithin{equation}{section}
\numberwithin{figure}{section}
\newlength{\lyxlabelwidth}      % auxiliary length 

%%%%%%%%%%%%%%%%%%%%%%%%%%%%%% User specified LaTeX commands.
\usepackage{fancyhdr}
\usepackage{lastpage}
\pagestyle{fancy}
\usepackage{enumitem}
\usepackage{ifthen}
%\everymath=\expandafter{\the\everymath\displaystyle}
%\DeclareMathSizes{9}{9}{9}{9}
\usepackage{multicol}

\usepackage{tikz}
\usetikzlibrary{patterns}
\usetikzlibrary{plotmarks}
\usepackage{pgfplots}
\pgfplotsset{compat=newest}

\usepackage{calc}
\usepackage[nomessages]{fp}

\usepackage{datetime}
% Added by lyx2lyx
\setlength{\parskip}{\smallskipamount}
\setlength{\parindent}{0pt}

\makeatother

\usepackage{babel}
\begin{document}
\renewcommand{\author}{Dr.\,James Ashe}

\newcommand{\classNum}{232}

\newcommand{\classSec}{B}

\newcommand{\examType}{Final Exam Review}

\renewcommand{\date}{Final date and time:  \formatdate{16}{5}{2013}~ 3-5pm}

%%First page's header%%%%%%%%%%%%%%%%%%%%%%
\fancypagestyle{firststyle} %%header settings for first pagr
{
	\fancyhead[L]{MTH \classNum{}(\classSec{}) \author{}\\
	\textbf{\examType{}} ~\date{}}
	\fancyhead[C]{}
	\fancyhead[R]{\textbf{}}
	\setlength{\headsep}{14pt}
}
%%%%%%%%%%%%%%%%%%%%%%%%%%%%%%%%%%%%%%%%%%

%%Next page's header%%%%%%%%%%%%%%%%%%%%%%%
\setlist{nolistsep}
\lhead{\classNum{}(\classSec{})} 
\chead{} 
\cfoot{\thepage{}~of~\pageref{LastPage}} 
\setlength{\headsep}{5pt} 
%%%%%%%%%%%%%%%%%%%%%%%%%%%%%%%%%%%%%%%%%%%

%%Various settings; see the preamble for other settings%%%%%
%\setenumerate[0]{leftmargin=12pt,itemindent=0pt,labelwidth=0pt}
\setlist{topsep=.5pt}
\renewcommand{\labelenumi}{\textbf{\arabic{enumi}.}}
\renewcommand{\labelenumii}{\textbf{\alph{enumii}.}}
\thispagestyle{firststyle}
%%%%%%%%%%%%%%%%%%%%%%%%%%%%%%%%%%%%%%%%%%

\newcommand{\xmax}{0}
\newcommand{\xmin}{-\xmax}
\newcommand{\xstep}{0}
\newcommand{\ymax}{0}
\newcommand{\ymin}{-\ymax}
\newcommand{\ystep}{0} 
\newcommand{\dspace}{\vspace{1cm}}

\small

\subsubsection*{Be able to recognize AND use the following formulas:}
\begin{itemize}
\item The average value of a function $f$ on the interval $[a,b]$ is given
by $\frac{1}{b-a}{\displaystyle \int_{a}^{b}f(x)\,dx}$
\item The \emph{length of the arc }on the graph of $y=f(x)$ from the point
$\left(a,f(a)\right)$ to the point $\left(b,f(b)\right)$ is ${\displaystyle \pi\int_{a}^{b}\sqrt{1+\left[f'(x)\right]^{2}}\,dx}$
\item The area of a region between $f(x)$ and $g(x)$ on $[a,b]$ is${\displaystyle \int_{a}^{b}\left(f(x)-g(x)\right)\,dx}$
\item The volume of the solid generated by revolving the region bound by
the $x$-axis and $f(x)$ from $[a,b]$ is $\pi{\displaystyle \int_{a}^{b}\left(f(x)\right)^{2}\,dx}$.
\item Integration by parts ${\displaystyle \int u\,dv=uv-{\displaystyle \int v\,du}}$.
\end{itemize}

\subsubsection*{Be able to recognize and use the following properties of integration:}
\begin{itemize}
\item FTC part 1: $F(x)={\displaystyle \int_{a}^{x}f(t)\,dt}$, this shows
that an indefinite integral can be reversed by differentiation (and
vice-versa up to a constant).

\begin{itemize}
\item [ex.]If $F(x)={\displaystyle \int_{0}^{x}t^{3}}\,dt$ then $\frac{d}{dx}F(x)={\displaystyle \int_{0}^{x}t^{3}}\,dt=x^{3}$,
and If $F(x)={\displaystyle \int_{0}^{x^{2}}t^{3}}\,dt$ then $\frac{d}{dx}F(x^{2})={\displaystyle \int_{0}^{x^{2}}t^{3}}\,dt=2x(x^{2})^{3}=2x^{7}$
\end{itemize}
\item FTC part 2: ${\displaystyle \int_{a}^{b}f(x)\,dx}=F(b)-F(a)$ where
$\frac{d}{dx}F(x)=f(x)$ on $[a,b]$.

\begin{itemize}
\item [ex.]${\displaystyle \int_{2}^{5}x^{2}\,dx}=\frac{x^{3}}{3}\bigr|_{2}^{5}=\frac{5^{3}}{3}-\frac{2^{3}}{3}=39$
\end{itemize}
\item If $\int_{a}^{b}f(x)\,dx=A$ then $\int_{b}^{a}f(x)\,dx=-A$
\item $\int_{a}^{b}f(x)+g(x)\,dx=\int_{a}^{b}f(x)\,dx+{\displaystyle \int_{a}^{b}}g(x)\,dx$.
\item $\int_{a}^{c}f(x)\,dx=\int_{a}^{b}f(x)\,dx+\int_{b}^{c}f(x)\,dx$
if $a\leq b\leq c$.
\item If $f(x)$ is an \emph{even }function, i.e., $f(x)=f(-x)$, then $\int_{-a}^{a}f(x)\,dx=2\int_{0}^{a}f(x)\,dx$.
\item If $f(x)$ is an \emph{odd }function, i.e., $f(x)=-f(x)$, then $\int_{-a}^{a}f(x)\,dx=0$.
\end{itemize}

\subsubsection*{Be able to recognize the techniques needed to integrate and be able
to use them:}
\begin{itemize}
\item substitution.
\item integration by parts
\item partial fractions
\item trigonometric powers
\item the trigonometric substitution using $\sqrt{a^{2}-x^{2}}=a\cos\theta$.
\end{itemize}

\subsubsection*{Be able to complete the following type of problems.}
\begin{itemize}
\item Initial value problems

\begin{itemize}
\item [ex.]$\frac{dy}{dx}=(x+2)^{.\nicefrac{3}{2}}$ and $y(0)=1$.
\end{itemize}
\item Definite integral of an absolute value.
\item Definite integral from a graph.\newpage\footnotesize
\end{itemize}

\subsection*{Exam 1}

\emph{Evaluate the following using the provided values.}\\
${\displaystyle \int_{0}^{\nicefrac{\pi}{2}}g(x)}\,dx=1$, ${\displaystyle \int_{\nicefrac{\pi}{2}}^{\pi}g(x)}\,dx=\mbox{\ensuremath{\pi}-1}$,
${\displaystyle \int_{\pi}^{\nicefrac{3\pi}{2}}g(x)}\,dx=\pi+1$,
${\displaystyle \int_{\nicefrac{3\pi}{2}}^{2\pi}g(x)}\,dx=2\pi-1$
\begin{enumerate}
\item ${\displaystyle \int_{0}^{\pi}g(x)}\,dx$
\end{enumerate}
\begin{enumerate}[resume]
\item ${\displaystyle \int_{\pi}^{0}\frac{g(x)}{\pi}}\,dx$
\item Suppose that $g(x)$ is \emph{even}, \\
find ${\displaystyle \int_{-\nicefrac{\pi}{2}}^{\nicefrac{\pi}{2}}g(x)}\,dx$.
\end{enumerate}
\emph{Use geometry or symmetry to evaluate the definite integral.
Recall that the area of a triangle is $A=\frac{1}{2}bh$.}
\begin{enumerate}[resume]
\item ${\displaystyle \int_{-2}^{2}\left(2-|x|\right)\,dx}$
\end{enumerate}
\emph{Evaluate the following.}
\begin{enumerate}[resume]
\item Let $a$ and $b$ be constants.

\begin{enumerate}[resume]
\item ${\displaystyle {\displaystyle \frac{d}{dx}\int_{a}^{x}3t^{2}\,dt}}$ 
\item ${\displaystyle {\displaystyle \frac{d}{dx}\int_{a}^{b}f(t)\,dt}}$
\end{enumerate}
\item $\int_{0}^{2}(x^{2}+1)\,dx$
\item $\int_{0}^{\pi}\frac{1}{2}\sin\frac{\theta}{2}\,d\theta$.
\item ${\displaystyle \int\frac{1-z^{4}}{z^{2}}}\,dz$
\item ${\displaystyle \int}x^{^{3}}(x^{4}+4)^{15}\,dx$
\item $\int\frac{1}{(5x-3)^{2}}\,dx$
\item ${\displaystyle \int\cos x\sin^{3}x\,dx}$
\end{enumerate}
\emph{Recall that $\sin x$ is an }\emph{\uline{odd}} \emph{function. }
\begin{enumerate}[resume]
\item ${\displaystyle {\displaystyle \int_{-5}^{5}}\sin^{3}x\,dx}$
\end{enumerate}
\emph{Solve the following word problems. }
\begin{enumerate}[resume]
\item An arch is 10 ft high and a 20 ft base. Its shape can be modeled by
$y=10-\frac{x^{2}}{10}$. Find the average height of the arch above
the ground. Note that the maximum height is at $x=0$.
\end{enumerate}

\subsection*{Exam 2}

\emph{Find the displacement over the given interval.}
\begin{enumerate}
\item $v(t)=\cos\left(\frac{\pi}{2}t\right)$ on $[0,3]$; $s(0)=0$.
\end{enumerate}
\emph{Determine the position function.}
\begin{enumerate}[resume]
\item An object is moving at a velocity of $(t+9)^{\nicefrac{1}{2}}$ with
an initial position of $s(0)=20$ at time $t$.
\end{enumerate}
\emph{Solve the problem.}
\begin{enumerate}[resume]
\item A projectile is launched vertically from the ground at $t=0$ where
$t$ is seconds, and its velocity in flight (in $\nicefrac{\mbox{m}}{\mbox{s}}$)
is given by $v(t)=20-10t$.

\begin{enumerate}
\item Find the position at time $t$ for $0\leq t\leq4$.
\item Find the total distance traveled from $t=0$ to $t=4$.
\end{enumerate}
\end{enumerate}
\emph{Use calculus to find the area of the region described.}
\begin{enumerate}[resume]
\item The region bounded by $y=x$ and $y=x^{2}-4x$.\\
\includegraphics[width=1.5cm]{D:/home/jim/JamesAshe/JCSU/Calc_2_MTH_232/images/x_and_x2-4x.png}
\item The region bounded by $y=\frac{1}{2}x^{2}$ and $y=-x^{2}+6$.\\
\includegraphics[width=1.5cm]{D:/home/jim/JamesAshe/JCSU/Calc_2_MTH_232/images/1-x_and_x-1.png}
\item The region in the first quadrant bounded by $-x+1$ and $x-1$ from
$0\leq x\leq2$.\\
\includegraphics[width=1.5cm]{D:/home/jim/JamesAshe/JCSU/Calc_2_MTH_232/images/1-x_and_x-1.png}
\end{enumerate}
\emph{Let $R$ be the region bounded by the following curves. Use
the disk method to find the volume of the solid generated when $R$
is revolved about the $x$-axis.}
\begin{enumerate}[resume]
\item $y=\sqrt{x}$, $x=0$, and $x=9$.
\item $y=\sqrt{\sin4x}$, $y=0$, $0\leq x\leq\frac{\pi}{4}$.
\end{enumerate}
\emph{Find the length of the curve.}
\begin{enumerate}[resume]
\item $y=2x^{\nicefrac{3}{2}}$ from $x=0$ to $x=\frac{5}{4}$.
\item $y=\left(4-x^{\nicefrac{2}{3}}\right)^{\nicefrac{3}{2}}$ from $x=0$
to $x=8$.
\end{enumerate}

\subsection*{Exam 3}

\emph{Perform the indicated operation.}
\begin{enumerate}[resume]
\item ${\displaystyle \int\left(10e^{5x}+\frac{1}{x}\right)\,dx}$
\item ${\displaystyle \int\frac{dx}{3x-5}}$
\end{enumerate}
\emph{Perform the indicated operation.}
\begin{enumerate}[resume]
\item ${\displaystyle \int}2xe^{2x}\,dx$
\item ${\displaystyle \int x\cos\left(8x\right)\,dx}$
\item ${\displaystyle \int10x\ln(2x)\,dx}$
\end{enumerate}

\subsection*{Exam 4}
\begin{enumerate}
\item ${\displaystyle \int}\cos^{3}2x\,dx$
\item ${\displaystyle \int}\sin^{2}x\,dx$
\item ${\displaystyle \int\sin^{3}x}\cos x\,dx$
\item ${\displaystyle \int\frac{dx}{\sqrt{25-x^{2}}}}$
\item ${\displaystyle \int_{0}^{7}\sqrt{49-x^{2}}}\,dx$. (Hint:${\displaystyle \int\cos^{2}x\,dx=\frac{x}{2}+\frac{\sin(2x)}{4}+C}$)
\item ${\displaystyle \int\frac{x}{x^{2}+x-12}\,dx}$
\item ${\displaystyle \int}_{0}^{\ln2}4xe^{x}\,dx$
\end{enumerate}

\end{document}
